% ch07_u13 —— 原书 p.356–p.366（第 7 章 7.14.3–7.15.17）
% 底本：A. Beilinson, V. Drinfeld, Quantization of Hitchin's Integrable System
% and Hecke Eigensheaves
% 承接：上单元（p.346–p.355）末句断在「以 $\omega_{(K\backslash Y)}$ 记典范」，
% 本单元正文自「丛 $\omega_{K\backslash Y}=\det\Omega_{K\backslash Y}$ 沿投影…的拉回」续起。

\SourcePage{356}
丛 $\omega_{K\backslash Y}=\det\Omega_{K\backslash Y}$ 沿投影 $Y\to K\backslash Y$ 的拉回（回忆 $K\backslash Y$ 是光滑栈）。这是分次线丛，位于次数 $\dim K\backslash Y$。

若 $K_1,K_2\subset G$ 是如上两个合理的群子概形，则有典范同构
\begin{equation}\tag{378}
\omega_{(K_1\backslash Y)}=\lambda_{(\mathfrak{k}_1,\mathfrak{k}_2)}\otimes\omega_{(K_2\backslash Y)}
\end{equation}
它满足显然的传递性。事实上，为定义（378），只需考虑 $K_2\subset K_1$ 的情形。把光滑投影 $K_2\backslash Y\to K_1\backslash Y$ 的相对切丛拉回到 $Y$，得 $(\mathfrak{k}_1/\mathfrak{k}_2)\otimes\mathcal{O}_Y$，由此得到（378）。传递性显然。

今 $C_Y^{\cdot}\in\mathcal{C}(Y)$ 是 Clifford 模，并附有一族同构 $\lambda_{\mathfrak{k}}(C_Y^{\cdot})\simeq\omega_{(K\backslash Y)}$，对每个合理的子群 $K\subset G$ 给出，且与（377）、（378）相容。这样的 $C_Y^{\cdot}$ 存在，并在唯一同构的意义下唯一。

$G$ 在 $Y$ 上的作用典范地提升为 $C_Y^{\cdot}$ 上的 $G$-作用，且与 $G$ 在 Clifford 算子 $\mathfrak{g}\oplus\mathfrak{g}^{*}$ 上的伴随作用相容。实际上，$G(\mathbb{C})$ 作用在构造 $C_Y^{\cdot}$ 所涉及的全部对象上，故它作用在 $C_Y^{\cdot}$ 上。要定义 $A$-点 $G(A)$ 在 $C_Y^{\cdot}\otimes A$ 上的作用，须用 $G$ 的 $A$-族合理群子概形，把 $Y\times\operatorname{Spec}A$ 上 Clifford 模 $C_Y^{\cdot}\otimes A$ 的特征性质写清楚。我们把它留给读者。

\emph{注.} 任取 $y\in Y$。$C_y^{\cdot}$ 在 $y$ 处的纤维 $C_y^{\cdot}$ 是不可约分次 $\mathrm{Cl}^{\cdot}_{\mathfrak{g}}$-模，可如下描述。% [Tutor Note: 原书 p.356 此处印作「$C_y^{\cdot}$ 在 $y$ 处的纤维 $C_y^{\cdot}$」，第二个下标当作 $Y$，底本如此，保留原样。]
考虑「作用」映射 $\mathfrak{g}\to\Theta_y$。其核 $\mathfrak{g}_y$（$y$ 的稳定化子）是 $\mathfrak{g}$ 中的 d-格。余核 $T$ 是有限维向量空间。令 $C_{y\mathfrak{g}_y}^{\cdot}$ 为 $C_y^{\cdot}$ 中 $\mathfrak{g}_y$-余不变量（关于 $\mathfrak{g}_y$ 的 Clifford 作用）构成的分次向量空间。则有典范等同 $C_{y\mathfrak{g}_y}^{\dim T}\simeq\det(T^{*})$，而 $C_y^{\cdot}$ 由这一规范化唯一确定。

\bdp{7.14.4} 设 $\mathcal{L}=\mathcal{L}_Y$ 是 $Y$ 上的线丛，带有提升 $Y$ 上 $G$-作用的 $G'$-作用；我们假定 $\mathbb{G}_m\subset G$ 以与标准相反的特征标作用在 $\mathcal{L}$ 上。

\SourcePage{357}
取 $V\in\mathcal{M}(\mathfrak{g})'$，于是 $V$ 是离散 $\mathfrak{g}'$-模，其中 $\mathbb{C}\subset\mathfrak{g}'$ 按标准特征标作用。此时张量积 $\mathcal{L}\otimes V$ 是 $\mathfrak{g}$-模，$C_Y^{\cdot}\otimes\mathcal{L}\otimes V$ 亦然（即 $Y$ 上的 $\mathfrak{g}$-作用提升为这些 $\mathcal{O}$-模上的连续 $\mathfrak{g}$-作用）。我们把 $\alpha\in\mathfrak{g}$ 在 $C_Y^{\cdot}\otimes\mathcal{L}\otimes V$ 上的作用记为 $\operatorname{Lie}_{\alpha}$。注意 $C_Y^{\cdot}\otimes\mathcal{L}\otimes V$ 还是 Clifford 模，且上述 $\mathfrak{g}$-作用与 Clifford 算子相容。照例把 $\alpha\in\mathfrak{g}$ 的 Clifford 作用记为 $i_{\alpha}$。于是与 7.13.26 一样，$C_Y^{\cdot}\otimes\mathcal{L}\otimes V$ 是分次 $(\Omega_{\mathfrak{g}},\mathfrak{g}_{\Omega})$-模。

下述命题与 7.13.27 类似，其证明亦然，我们留给读者。

\bdp{7.14.5} \emph{命题.} 存在唯一层态射
\[
d:C_Y^{\cdot}\otimes\mathcal{L}\otimes V\to C_Y^{\cdot+1}\otimes\mathcal{L}\otimes V,
\]
使得对任意 $\alpha\in\mathfrak{g}$ 有 $\operatorname{Lie}_{\alpha}=di_{\alpha}+i_{\alpha}d$。这个 $d$ 是一阶微分算子，$d^2=0$，且 $C_{\mathcal{L}}(V):=(C_Y^{\cdot}\otimes\mathcal{L}\otimes V,d)$ 是 DG $(\Omega_{\mathfrak{g}},\mathfrak{g}_{\Omega})$-模。

\emph{注.} 也可由 7.13.27 直接推出 7.14.5。取任意如 7.14.3 的 $K$。则 $C_Y^{\cdot}\otimes\omega^{\ast}_{(K\backslash Y)}$ 是「常」Clifford 模：它典范同构于 $C^{\cdot}\otimes\mathcal{O}_Y$，其中 $C^{\cdot}$ 是某个不可约 Clifford 模。$C^{\cdot}$ 上的 $\mathfrak{g}^{\flat}$-作用与 $C_Y^{\cdot}$ 上的 $\mathfrak{g}^{\flat}$-作用给出 $\omega_{(K\backslash Y)}=\operatorname{Hom}(C_Y^{\cdot},C^{\cdot}\otimes\mathcal{O}_Y)$ 上的 $\mathfrak{g}^{\flat}$-作用，它提升 $Y$ 上的 $\mathfrak{g}$-作用。于是 $\mathfrak{g}^{\flat}$ 作用在 $\omega_{(K\backslash Y)}\otimes\mathcal{L}\otimes V$ 上，而 7.14.5 中的 $d$ 与 7.13.27 中关于 $C^{\cdot}\otimes(\omega_{(K\backslash Y)}\otimes\mathcal{L}\otimes V)$ 的 $d$ 一致。

\bdp{7.14.6} 这样我们定义了 $Y$ 上的 Ω-复形 $C_{\mathcal{L}}(V)$。当 $V$ 是 $\mathcal{M}(\mathfrak{g})'$ 中的复形时，用显然的方式延拓此定义。今设 $K$ 如 7.14.1。对 Harish-Chandra 复形 $V\in C(\mathfrak{k}_{\Omega}\times\mathfrak{g},K)'$，Ω-复形 $C_{\mathcal{L}}(V)$ 是 $K_{\Omega}$-等变的。实际上，$K$ 依其在 $C_Y^{\cdot}$、$\mathcal{L}$、$V$ 上的作用作用在 $C_{\mathcal{L}}(V)$ 上，而算子 $i_{\xi}$（$\xi\in\mathfrak{k}$）等于 $C_Y^{\cdot}$ 的相应 Clifford 算子与 $V$ 上 $\mathfrak{k}_{\Omega}^{-1}$-作用的算子之和。

\SourcePage{358}
置 $\Delta_{\Omega\mathcal{L}}(V):=C_{\mathcal{L}}(V)[\dim(K\backslash Y)]$。我们已定义 DG 函子
\begin{equation}\tag{379}
\Delta_{\Omega}=\Delta_{\Omega\mathcal{L}}:C(\mathfrak{k}_{\Omega}\times\mathfrak{g},K)'\longrightarrow C(K\backslash Y,\Omega)
\end{equation}

\bdp{7.14.7} \emph{注.} Ω-复形 $\Delta_{\Omega}(V)$ 带有典范滤过 $\Delta_{\Omega}(V)_{\cdot}$，其中 $\Delta_{\Omega}(V)_a$ 由这样的截面组成：$\mathfrak{k}^{\perp}\subset\mathfrak{g}^{*}$ 中任意 $a+1$ 个 Clifford 算子的乘积作用在它上面为零（见 7.13.29）。% [Tutor Note: 原书 p.358 此处印作「filtrartion」，应为「filtration」，底本如此，保留原样。]
由（376）有 $K_{\Omega}$-等变 Ω-复形的典范同构
\begin{equation}\tag{380}
\operatorname{gr}_a\Delta_{\Omega}(V)\xrightarrow{\ \sim\ }C_{\mathfrak{k}}(\omega_{(K\backslash Y)}\otimes\mathcal{L}\otimes V\otimes\Lambda^{a}(\mathfrak{g}/\mathfrak{k}))[a]
\end{equation}

\bdp{7.14.8} \emph{引理.} (i) 函子 $\Delta_{\Omega}$ 把拟同构送到 D-拟同构，故给出三角函子
\begin{equation}\tag{381}
L\Delta=L\Delta_{\mathcal{L}}:D(\mathfrak{g},K)'\to D(K\backslash Y)
\end{equation}
(ii) 函子 $L\Delta$ 右 t-正合，相应的右正合函子 $\Delta=\Delta_{\mathcal{L}}:\mathcal{M}((\mathfrak{g},K)'\to\mathcal{M}^{\ell}(K\backslash Y)$ 为
\begin{equation}\tag{382}
\Delta_{\mathcal{L}}(V)_Y=(\mathcal{D}_Y\otimes\mathcal{L})\widehat{\otimes}_{U(\mathfrak{g})}V=\mathcal{L}^{*}\otimes\mathcal{D}_{Y,\mathcal{L}}\widehat{\otimes}_{U(\mathfrak{g}')}V
\end{equation}
这里 $\mathcal{D}_Y$ 是 $Y$ 上微分算子的拓扑代数（见 1.2.6），$\mathcal{D}_{Y,\mathcal{L}}:=\mathcal{L}\otimes\mathcal{D}_Y\otimes\mathcal{L}^{*}$ 是相应的 $\mathcal{L}$-扭代数。% [Tutor Note: 原书 p.358 此处印作「$\mathcal{M}((\mathfrak{g},K)'$」，左括号多了一个，底本如此，保留原样。]

\emph{证明.} (i) 我们的命题是局部的，故必要时缩小 $K$，可设 $K$ 在 $Y$ 上的作用自由。把 $\Delta_{\Omega}(V)$ 视为 $K\backslash Y$ 上的滤过 Ω-复形。对 $K$-模 $P$，记 $P^{\sim}$ 为其 $Y$-扭，即 $K\backslash Y$ 上的 $\mathcal{O}$-模。投影 $C_{\mathfrak{k}}\to C_{\mathfrak{k}}/C_{\mathfrak{k}}^{\ge1}$ 按（380）给出典范同构
\begin{equation}\tag{383}
\operatorname{gr}_a\Delta_{\Omega}(V)_{K\backslash Y}=\omega_{K\backslash Y}\otimes\mathcal{L}_{K\backslash Y}\otimes V^{\sim}\otimes\Lambda^{a}(\mathfrak{g}/\mathfrak{k})^{\sim}[a]
\end{equation}
右端是 $\mathcal{O}$-复形，故 $V$ 之间的拟同构给出 $\Delta_{\Omega}(V)$ 之间的（滤过）D-拟同构。

(ii) 同上可设 $K$ 的作用自由。对 $V\in\mathcal{M}(\mathfrak{g},K)'$，可把（383）改写为同构 $\Delta_{\Omega}(V)^{a}_{K\backslash Y}=\omega_{K\backslash Y}\otimes\mathcal{L}_{K\backslash Y}\otimes V^{\sim}\otimes\Lambda^{-a}(\mathfrak{g}/\mathfrak{k})^{\sim}$。这说明 $\Delta_{\Omega}$ 右 t-正合。\SourcePage{359}%
$\Delta_{\Omega}(V)_{K\backslash Y}$ 中的微分如下。% [Tutor Note: 原书 p.359 此处印作「The differential in $\Delta_{\Omega}(V)_{K\backslash Y}$ as follows.」，缺谓语（疑脱 is），底本如此，译文按「如下」处理。]
$Y$ 上的 $\mathfrak{g}$-作用在 $(\mathfrak{g}/\mathfrak{k})^{\sim}$ 上定义了 $K\backslash Y$ 上 Lie 代数胚的结构。$\mathcal{L}_Y\otimes V$ 上的 $\mathfrak{g}$-作用在 $\mathcal{L}_{K\backslash Y}\otimes V^{\sim}$ 上定义了左 $(\mathfrak{g}/\mathfrak{k})^{\sim}$-模结构，故 $\omega_{K\backslash Y}\otimes\mathcal{L}_{K\backslash Y}\otimes V^{\sim}$ 是右 $(\mathfrak{g}/\mathfrak{k})^{\sim}$-模。于是 $\Delta_{\Omega}(V)_{K\backslash Y}$ 是 $(\mathfrak{g}/\mathfrak{k})^{\sim}$ 的、系数在 $\omega_{K\backslash Y}\otimes\mathcal{L}_{K\backslash Y}\otimes V^{\sim}$ 中的 Chevalley 同调复形。$K\backslash Y$ 上的右 D-模 $H^0_D(L\Delta(V))$ 为 $(\omega_{K\backslash Y}\otimes\mathcal{L}_{K\backslash Y}\otimes V^{\sim})\otimes_{(\mathfrak{g}/\mathfrak{k})^{\sim}}\mathcal{D}_{K\backslash Y}$；相应的左 D-模是 $\mathcal{D}_{K\backslash Y}\otimes_{(\mathfrak{g}/\mathfrak{k})^{\sim}}(\mathcal{L}_{K\backslash Y}\otimes V^{\sim})$。把这一同构提升到 $Y$ 上即得（382）。

\bdp{7.14.9} \emph{例.} 我们来计算 $L\Delta(Vac')$。嵌入 $\mathbb{C}\to Vac'$ 给出 $Y$ 上 Ω-复形的嵌入 $C_{\mathfrak{k}}(\omega_{(K\backslash Y)}\otimes\mathcal{L}_Y)\to\Delta_{\Omega\mathcal{L}}(Vac')_0$。相应的 $K_{\Omega}$-等变 Ω-复形态射 $C_{\mathfrak{k}}(\omega_{(K\backslash Y)}\otimes\mathcal{L}_Y)\to\Delta_{\Omega\mathcal{L}}(Vac')$ 是 D-拟同构，这一验证留给读者。今左端是 $K\backslash Y$ 上的 Ω-复形 $\Omega(\mathcal{D}_{K\backslash Y}\otimes\mathcal{L}_{K\backslash Y}$（见 7.3.3）。% [Tutor Note: 原书 p.359 此处括号不配对，右括号缺失，底本如此，保留原样。]
因此若 $K\backslash Y$ 是良栈，则
\[
L\Delta(Vac')=\Delta(Vac')=\mathcal{D}_{K\backslash Y}\otimes\mathcal{L}_{K\backslash Y}.
\]

\emph{注.} 由于 $\operatorname{End}Vac'$ 反同构于 1.2.5 中的代数 $D'_{(\mathfrak{g},K)}$（另见 1.2.2），我们有 $D'_{(\mathfrak{g},K)}$ 在 $\Delta(Vac')=\mathcal{D}_{K\backslash Y}\otimes\mathcal{L}_{K\backslash Y}$ 上的右作用，即从 $D'_{(\mathfrak{g},K)}$ 到扭微分算子环 $\Gamma(K\backslash Y,D'_{K\backslash Y})$ 的同态。这就是 1.2.5 中的同态 $h$（另见 1.2.3 与 1.2.4）。

\bdp{7.14.10} \emph{命题.} 函子 $L\Delta:D(\mathfrak{g},K)'\to D(K\backslash Y)$ 是 H-模的态射。

\emph{证明.} 7.8.8 的构造与论证以显然的方式移植到我们的无穷维情形。7.9 的无穷维版本是直接的。

\SourcePage{360}
\bdsec{7.15. 仿射旗空间是 \texorpdfstring{$\mathcal{D}$}{D}-仿射的}

本节说明：水平低于临界值的仿射李代数表示与仿射旗空间上的 $\mathcal{D}$-模之间的关系，与通常的有限维情形相同。

\bdp{7.15.1} 下文恒设 $K=\mathbb{C}((t))$，$O=\mathbb{C}[[t]]$。设 $\mathfrak{g}$ 是单（有限维）李代数\footnote{当 $\mathfrak{g}$ 是任意 reductive 李代数时，此推广是直接的。}，$G$ 是相应的单连通单群。我们有群归纳概形 $G(K)$ 及其群子概形 $G(O)$（见 7.11.2(iv)）。$G(K)$ 在 Tate 向量空间 $\operatorname{Lie}G(K)=\mathfrak{g}(K)$ 上的伴随作用给出 $G(K)$ 被 $\mathbb{G}_m$ 的中心扩张 $G(K)^{\flat}$（见 ??）。其李代数是 $\mathfrak{g}(K)$ 的中心扩张 $\mathfrak{g}(K)^{\flat}$，由余圈 $\phi,\psi\mapsto\operatorname{Res}(d\phi,\psi)$ 定义，其中 $(a,b):=\operatorname{Tr}(\operatorname{ad}_a\cdot\operatorname{ad}_b)$（见 ??）。令 $G(O)^{\flat}\subset G(K)^{\flat}$ 为 $G(O)$ 的原像。$G(O)$ 的伴随作用保持 c-格 $\mathfrak{g}(O)\subset\mathfrak{g}(K)$，故有典范等同 $s:G(O)^{\flat}\xrightarrow{\sim}G(O)\times\mathbb{G}_m$\footnote{由于 $G$ 单，分裂 $G(O)\to G(O)^{\flat}$ 唯一。}。% [Tutor Note: 原书 p.360 两处交叉引用均印作「??」（引用未解析），底本如此，保留原样。]

设 $N\subset B\subset G$，其中 $B$ 是 Borel 子群，$N$ 是其幂幺根，故 $H=B/N$ 是 $G$ 的嘉当群。令 $N^{+},B^{+}$ 为 $N,B$ 在显然投影 $G(O)\to G$ 下的原像，故 $B^{+}/N^{+}=H$，$G(O)/B^{+}=G/B$。令 $B^{\dagger}\subset G(K)^{\flat}$ 为 $B^{+}$ 的原像。存在唯一截面 $N^{+}\to G(K)$；置 $H^{\flat}:=B^{\dagger}/N^{+}$，$\mathfrak{h}^{\flat}=\operatorname{Lie}H^{\flat}$。截面 $s$ 给出同构 $B^{+}\times\mathbb{G}_m\xrightarrow{\sim}B^{\dagger}$，从而有同构 $H\times\mathbb{G}_m\xrightarrow{\sim}H^{\flat}$、$\mathfrak{h}\times\mathbb{C}\xrightarrow{\sim}\mathfrak{h}^{\flat}$。

置 $X:=G(K)/B^{+}=G(K)^{\flat}/B^{\dagger}$（层的商，取平坦或 Zariski 拓扑都一样，这一点由 4.5.1 推出）。$X$ 称为\emph{仿射旗空间}。它是既约、连通、归纳投射、形式光滑的归纳概形\footnote{$X$ 光滑地纤维化到仿射 Grassmann 流形 $G(K)/G(O)$ 上，见 4.5.1。}。置 $X^{\dagger}:=G(K)^{\flat}/N^{+}$：这是 $X$ 上的左 $H^{\flat}$-torsor（主齐性空间），作用为 $h^{\flat}\cdot x^{\dagger}=x^{\dagger}h^{\flat-1}$。它带有显然的 $G(K)^{\flat}$-作用。记投影 $X^{\dagger}\to X$ 为 $p$。

\SourcePage{361}
\bdp{7.15.2} 令 $\mathcal{M}^{\dagger}(X)$ 为 $X^{\dagger}$ 上弱 $H^{\flat}$-等变 $\mathcal{D}$-模的范畴（见 7.11.11）。这是 Abel 范畴。对 $M\in\mathcal{M}^{\dagger}(X)$ 置 $M_X:=(p_{\cdot}M)^{H^{\flat}}\in\mathcal{M}(X,\mathcal{O})$。函子 $\mathcal{M}^{\dagger}(X)\to\mathcal{M}(X,\mathcal{O})$，$M\mapsto M_X$，正合且忠实。

置 $\mathcal{D}^{\dagger}:=(p_{\cdot}\mathcal{D}^{\dagger}_X)^{H^{\flat}}$。这是 $X$ 上的 Diff-代数。映射
\begin{equation}\tag{384}
\mathfrak{h}^{\flat}\to\Gamma(X,\mathcal{D}^{\dagger})=\Gamma(X^{\dagger},\mathcal{D}^{\dagger}_{X^{\dagger}})^{H^{\flat}}
\end{equation}
等于沿 $p$ 的纤维的左作用的相反数，其像落在 $\mathcal{D}^{\dagger}$ 的中心内。实际上，$\mathcal{D}^{\dagger}$ 是 $\operatorname{Sym}(\mathfrak{h}^{\flat})$-族的 tdo（见 7.11.11(b)）。

注意 $\mathcal{D}^{\dagger}$ 以显然的方式（从右）作用在如上任何 $M_X$ 上，故有函子
\begin{equation}\tag{385}
\mathcal{M}^{\dagger}(X)\to\mathcal{M}(X,\mathcal{D}^{\dagger}).
\end{equation}
有（见 7.11.11 中的注 (ii)）：

\bdp{7.15.3} \emph{引理.} 函子（385）是范畴等价。\hfill$\Box$

\bdp{7.15.4} 对 $\chi=(\chi_0,c)\in\mathfrak{h}^{\flat*}=\mathfrak{h}^{*}\times\mathbb{C}$，记 $\mathcal{D}^{\chi}$ 为我们族 $\mathcal{D}^{\dagger}$ 中相应的 tdo。于是 $\mathcal{D}^{(0,0)}=\mathcal{D}_X$。置 $\mathcal{M}^{\chi}(X):=\mathcal{M}(X,\mathcal{D}^{\chi})\subset\mathcal{M}(X,\mathcal{D}^{\dagger})$。考虑拓扑代数 $\Gamma\mathcal{D}^{\chi}=\Gamma(X,\mathcal{D}^{\chi})$（见 7.11.9、7.11.10）。我们有函子
\begin{equation}\tag{386}
\Gamma:\mathcal{M}^{\chi}(X)\to\mathcal{M}^{r}(\Gamma\mathcal{D}^{\chi})
\end{equation}
其中 $\mathcal{M}^{r}(\Gamma\mathcal{D}^{\chi})$ 是离散右 $\Gamma\mathcal{D}^{\chi}$-模的范畴，$\Gamma M:=\Gamma(X,M)$。

$\mathfrak{g}(K)^{\flat}$ 在 $X^{\dagger}$ 上的作用给出连续态射 $\mathfrak{g}(K)^{\flat}\to\Gamma(X,\mathcal{D}^{\dagger})$。相应的态射 $\mathfrak{g}(K)^{\flat}\to\Gamma\mathcal{D}^{\chi}$ 把 $1^{\flat}\in\mathfrak{g}(K)^{\flat}$ 送到 $-c$。

\bdp{7.15.5} 若 Verma $\mathfrak{g}(K)^{\flat}$-模 $M(\chi)$ 不可约，则称 $\chi$ 是\emph{反占优}的。据 [KK] 3.1，这等价于以下三个条件：
\begin{enumerate}[label=(\roman*),leftmargin=2.2em]
\item 有 $c\ne-1/2$。
\SourcePage{362}%
\item 对 $\mathfrak{g}$ 的任意正余根 $h_{\alpha}\in\mathfrak{h}$ 有 $(\chi_0+\rho_0)(h_{\alpha})\ne1,2,\ldots$。
\item 对如上任意 $h_{\alpha}$ 及任意整数 $n>0$ 有
\[
\pm(\chi_0+\rho_0)(h_{\alpha})+2n\frac{c+1/2}{(\alpha,\alpha)}\ne1,2,\ldots
\]
\end{enumerate}

这里 $\rho_0\in\mathfrak{h}^{*}$ 是 $\mathfrak{g}$ 的正根之和的一半，$(\cdot,\cdot)$ 是 $\mathfrak{h}^{*}$ 上与 $\mathfrak{h}$ 上 $(\cdot,\cdot)$ 相对应的标量积（见 7.15.1）。

\emph{注.} 为从 [KK] 3.1 推出上述论断，只需注意 $\mathfrak{g}(K)^{\flat}$ 的「实」正余根是 $h_{\alpha}$ 与 $\pm h_{\alpha}+2n(\alpha,\alpha)^{-1}\cdot1^{\flat}$，其中 $h_{\alpha}$、$n$ 如上；而 [KK] 中的权 $\rho$ 由下式给出。

置 $\rho:=(\rho_0,1/2)\in\mathfrak{h}^{\flat*}$。若 $\chi+\rho$ 在仿射 Weyl 群 $W_{\mathrm{aff}}$ 中的稳定化子平凡，则称 $\chi$ 是\emph{正则}的\footnote{回忆 $W_{\mathrm{aff}}$ 在 $\mathfrak{h}^{\flat*}$ 上的作用来自 $G(K)$ 在 $\mathfrak{g}(K)^{\flat}$ 上的伴随作用。}。

\bdp{7.15.6} \emph{定理.} 设 $\chi$ 反占优且正则。则（386）是范畴等价。

我们在下文 7.15.8--?? 中证明 7.15.6。

\bdp{7.15.7} \emph{注.} (i) 令 $\mathcal{M}^{c}(\mathfrak{g}(K))$ 为离散 $\mathfrak{g}(K)^{\flat}$-模的范畴，其中 $1^{\flat}$ 以 $c$ 为乘子作用。令
\begin{equation}\tag{387}
\Gamma:\mathcal{M}^{\chi}(X)\to\mathcal{M}^{c}(\mathfrak{g}(K))
\end{equation}
为（386）与显然的「限制」函子 $\mathcal{M}^{r}(\Gamma\mathcal{D}^{\chi})\to\mathcal{M}^{c}(\mathfrak{g}(K))$ 的合成。据 7.15.6，此函子正合且忠实。

(ii) 人们可以期望 $\mathfrak{g}(K)^{\flat}$ 在 $\Gamma\mathcal{D}^{\chi}$ 中生成稠密子代数\footnote{这等价于：对 $M\in\mathcal{M}^{\chi}(X)$，$\Gamma M$ 的任何 $\mathfrak{g}(K)^{\flat}$-子模都来自 $M$ 的 $\mathcal{D}^{\chi}$-子模。}。换言之，$\Gamma\mathcal{D}^{\chi\circ}$ 是水平 $c$ 的包络代数 $\bar U^{c}=\bar U^{c}\mathfrak{g}(K)$ 按某种拓扑的完备化。% [Tutor Note: 原书 p.362 此处上标印作 $\chi\circ$（$\chi$ 后带一小圆圈），同段前文印作 $\Gamma\mathcal{D}^{\chi}$，底本如此，保留原样。]
能否显式确定这一拓扑？

注意，在有限维情形（见 [BB81] 或 [Kas]），人们通常由相应的「经典」论断推出\SourcePage{363}%
（利用 Kostant 的正规性定理）。这个「经典」论断（它说 $\mathfrak{g}(K)\hookrightarrow\Gamma(X,\Theta_X)$ 在 $\bigoplus_{n\ge0}\Gamma(X,\Theta_X^{\otimes n})$ 中生成稠密子代数）对仿射旗情形是\emph{假的}（例如，映射 $\mathfrak{g}(K)\to\Gamma(X,\Theta_X)$ 不满）。

如同 [BB81] 或 [Kas]，容易看出 7.15.6 由下述论断推出：

\bdp{7.15.8} \emph{定理.} (i) 若 $\chi$ 反占优，则对任意 $M\in\mathcal{M}^{\chi}(X)$ 及任意 $r>0$ 有 $H^{r}(X,M)=0$\footnote{这里 $H^{r}(X,M):=\varinjlim H^{r}(Y,M_{(Y)})$；我们用 7.11.4 的记号。}。

(ii) 若此外 $\chi$ 正则且 $M\ne0$，则 $\Gamma M\ne0$。

\emph{注.} 7.15.8(i) 的证明与相应的有限维论断的证明十分相似（见 [BB81] 或 [Kas]）。若能找到与 [BB81]（用平移函子）类似的方法来证明 7.15.8(ii)，那将很好，因为它有助于理解 7.15.7(ii)。

\bdp{7.15.9} 我们开始证明 7.15.8(i)。设 $\psi=(\psi_0,b)$ 是 $H^{\flat}$ 的特征标，$\mathcal{L}=\mathcal{L}^{\psi}$ 是相应的 $X$ 上 $G(K)^{\flat}$-等变线丛（由 $X^{\dagger}$ 定义）。假定 $\mathcal{L}$ 丰沛。这等价\footnote{见 7.15.5 中的注。}于 $\psi$ 的下述性质：对 $\mathfrak{g}$ 的任意正余根 $h_{\alpha}$ 有 $\frac{2b}{(\alpha,\alpha)}<\psi_0(h_{\alpha})<0$。

记 $V$ 为 pro-有限维向量空间 $\Gamma(X,\mathcal{L})$ 的对偶。% [Tutor Note: 原书 p.363 此处印作「Denote by $V$ be the dual…」，多出一 be，底本如此，保留原意。]
它以显然的方式成为 $G(K)^{\flat}$-模，从而是水平 $-b$ 的可积 $\mathfrak{g}(K)^{\flat}$-模\footnote{据 Borel--Weil 定理的一个变体（例如见 [?]），$V$ 是不可约 $\mathfrak{g}(K)^{\flat}$-模。}。考虑 $V\widehat{\otimes}\mathcal{L}$ 的典范截面；这是 $\mathcal{O}^{\flat}$-模的 $G(K)^{\flat}$-等变态射 $\mathcal{O}_X\to V\widehat{\otimes}\mathcal{L}$。把它与 $M$ 张量，得 $\mathcal{O}^{!}$-模的态射
\begin{equation}\tag{388}
i:M\to V\otimes\mathcal{L}\otimes M
\end{equation}
它与 $\mathfrak{g}(K)^{\flat}$ 的作用交换。% [Tutor Note: 原书 p.363 此两处 $\mathcal{O}$ 的上标在扫描件中字形不一（前一处似 $\flat$，后一处为细竖笔），据 7.15.10--7.15.11 的 $\mathcal{O}^{!}$-模与 !-层记法，统一作 $\mathcal{O}^{!}$。]

\SourcePage{364}
\bdp{7.15.10} 下面我们考虑 $X$ 上向量空间的 \emph{!-层}。这种对象 $F$ 是一条规则：对闭子概形 $Y\subset X$ 指定 $Y$ 的 Zariski 拓扑上的层 $F_{(Y)}$，并对 $Y\subset Y'$ 指定等同 $i^{!}_{YY'}F_{(Y')}=F_{(Y)}$\footnote{这里 $i^{!}_{YY'}F_{(Y')}:=$ 在 $Y$ 上（集合论意义下）支撑的截面构成的子层。}，它们满足显然的传递性（对照 7.11.4 的注 (i)）。注意 !-层构成 Abel 范畴。它把 $Y$ 上的层范畴作为对子商与扩张封闭的全子范畴包含在内。$X$ 上任何 $\mathcal{O}^{!}$-模 $M$ 给出 $X$ 上的 !-层 $\varinjlim M_{(Y)}$（于是 $Y$ 上相应的层是 $M_{(Y^{\wedge})}$\footnote{记号见 7.11.4。}）；为简便起见，我们仍把它记为 $M$。我们还将考虑 $\mathfrak{g}(K)^{\flat}$-模的 !-层，即配有 $\mathfrak{g}(K)^{\flat}$-作用的向量空间 !-层，使得在每个 $F_{(Y)}$ 上的作用以显然的意义离散。任何带有 $\mathfrak{g}(K)^{\flat}$-作用的 $\mathcal{O}^{!}$-模都可视为 $\mathfrak{g}(K)^{\flat}$-模的 !-层。

\bdp{7.15.11} \emph{命题.} 作为 $\mathfrak{g}(K)^{\flat}$-模的 !-层之间的态射，（388）是直和项嵌入。

\bdp{7.15.12} \emph{7.15.8(i) 的证明.} 任取 $\alpha\in H^{r}(X,M)=\varinjlim H^{r}(X_{(Y)},M_{(Y)})$。它来自某个闭子概形 $Y\subset X$ 及 $\mathcal{O}$-凝聚子模 $F\subset M_{(Y)}$。选取如上丰沛的 $\mathcal{L}$，使 $H^{r}(Y,\mathcal{L}\otimes F)=0$。由于 $i(\alpha)$ 属于 $H^{r}(Y,V\otimes\mathcal{L}\otimes F)$ 的像，它为零。由 7.15.11 得证。\hfill$\Box$

\bdp{7.15.13} \emph{7.15.11 的证明.} 我们将定义 $V\otimes\mathcal{L}\otimes M$ 的自同态 $A$，使得
\begin{equation}\tag{389}
\operatorname{Ker}A=M,\quad V\otimes\mathcal{L}\otimes M=\operatorname{Ker}A\oplus\operatorname{Im}A.
\end{equation}
这就证明了 7.15.11。

令 $\bar U:=\bar U\mathfrak{g}(K)^{\flat}$ 为 $\mathfrak{g}(K)^{\flat}$ 的通常完备化包络代数。考虑由公式（85）定义的 Sugawara 元 $\widetilde{\mathfrak{L}}_0\in\bar U$。对任意 $ft^{r}\in\mathfrak{g}((t))\subset\bar U$ 有 $[\widetilde{\mathfrak{L}}_0,ft^{r}]=(1^{\flat}+1/2)rft^{r}$（见（87））。对任意\SourcePage{365}%
$N\in\mathcal{M}^{c}(\mathfrak{g}(K))$（其中 $e\ne-1/2$），考虑作用在 $N$ 上的算子 $\Delta_N:=(e+1/2)^{-1}\widetilde{\mathfrak{L}}_0$。若又有 $e-b\ne-1/2$，则置
\begin{equation}\tag{390}
A_{V,N}:=\Delta_{V\otimes N}-\Delta_V\otimes\operatorname{id}_N-\operatorname{id}_V\otimes\Delta_N\in\operatorname{End}(V\otimes N).
\end{equation}
此算子与 $\mathfrak{g}(K)^{\flat}$ 的作用交换。

我们把这一构造用于 $\mathfrak{g}(K)^{\flat}$-模的 !-层 $N:=\mathcal{L}\otimes M$（于是 $e=b+c$，水平的条件满足）。置
\begin{equation}\tag{391}
A:=A_{V,\mathcal{L}\otimes M}\in\operatorname{End}(V\otimes\mathcal{L}\otimes M).
\end{equation}
我们来证明 $A$ 满足（389）。\hfill$\Box$

\bdp{7.15.14} 现在转向 7.15.8(ii)。它是下述命题的直接推论，该命题特别说明如何用 $\Gamma M$ 计算 $M$ 的纤维。先从记号说起。

考察 $X$ 按 $N^{+}$-轨道（Schubert 胞腔）的分层。胞腔由仿射 Weyl 群 $W_{\mathrm{aff}}$ 的元素标号。对 $w\in W_{\mathrm{aff}}$，相应的胞腔为 $i_w:Y_w\hookrightarrow X$，其维数为 $l(w)$。$H^{\flat}$-torsor（主齐性空间）$X^{\dagger}$ 到 $Y_w$ 的限制平凡\footnote{在 $X^{\dagger}$ 中、位于 $Y_w$ 上方的任何 $N^{+}$-轨道都给出一个截面。}。由于 $Y_w$ 上任何可逆函数都是常数，平凡化在相差常数平移的意义下唯一。因此 tdo $\mathcal{D}^{\chi}$ 拉回到 $Y_w$ 上是典范平凡化的。

令 $M$ 是导出范畴 $D(X,\mathcal{D}^{\chi})$ 的任意对象\footnote{其定义与 7.11.14 在非扭情形给出的定义类似。}。对任意 $w\in W_{\mathrm{aff}}$，我们有（非扭的，正如刚才所说明）$\mathcal{D}$-复形 $i_w^{!}M\in D(Y_w)$。

我们想计算李代数（连续）上同调 $H^{a}(\mathfrak{n}^{+},\Gamma M)$（注意，由于 7.15.8(i)，$\Gamma=\mathrm{R}\Gamma$）。由于 $\mathfrak{h}^{\flat}=\mathfrak{b}^{\dagger}/\mathfrak{n}^{+}$，这些是 $\mathfrak{h}^{\flat}$-模。我们假定 $\chi$ 正则。

\bdp{7.15.15} \emph{命题.} 存在典范同构
\[
H^{a}(\mathfrak{n}^{+},\Gamma M)\simeq\bigoplus_{w\in W_{\mathrm{aff}}}H^{a-l(w)}_{DR}(Y_w,i_w^{!}M),
\]
\SourcePage{366}%
使得 $\mathfrak{h}^{\flat}$ 在第 $w$ 个直和项上以 $w(\chi)$ 为乘子作用\footnote{回忆 $G(K)$ 在 $\mathfrak{g}(K)^{\flat}$ 上的伴随作用给出 $\mathfrak{h}^{\flat}$ 上的 $W_{\mathrm{aff}}$-作用。}。

\bdp{7.15.16} \emph{7.15.8(ii) 的证明.} 由于 $\Gamma$ 正合，可设 $M$ 紧支且有限生成。令 $Y\subset X$ 是 $M$ 的（约化）支撑中的光滑 Zariski 开子集。则 $M_{(Y)}$ 是光滑概形 $Y$ 上的凝聚 $\mathcal{D}$-模。故再缩小 $Y$，可设 $M_{(Y)}$ 是自由 $\mathcal{O}_Y$-模。今对任意 $x\in Y$ 有 $H^{\cdot}i_x^{!}M\ne0$。平移 $M$ 可设 $x=Y_1$。由 7.15.15 有 $H^{\cdot}(\mathfrak{n}^{+},\Gamma M)\ne0$，故 $\Gamma M\ne0$。\hfill$\Box$

\bdp{7.15.17} \emph{7.15.15 的证明.} 可设 $M=i_{w*}N$，其中 $N\in D(Y_w)$。确实，任何 $M\in D(X,\mathcal{D}^{\chi})$ 带有典范滤过，其 $\operatorname{gr}_iM=\bigoplus_{l(w)=i}i_{w*}i_w^{!}M$。今 7.15.15 对 $M$ 的同构来自对 $i_{w*}i_w^{!}M$ 的相应同构，连同 $\mathfrak{h}^{\flat}$ 作用的谱分解。这里我们用到 $\chi$ 正则的假设；论证的其余部分只需 $\chi$ 反占优。

先考虑 $M=\delta$ 的情形，此时 $\Gamma\delta$ 是 7.15.5 中的 Verma 模（见 7.15.7(iii)）。% [Tutor Note: 原书 p.366 引「7.15.7(iii)」，但 7.15.7 只有 (i)(ii)，底本如此，保留原样。]
这个 Verma 模是秩 1 的余自由 $N^{+}$-模（它由任何不消去真空向量的泛函 $\nu$ 余自由生成）\footnote{$\nu$ 的核不含非平凡 $\mathfrak{n}^{+}$-子模（否则，由于 $\mathfrak{n}^{+}$ 幂零，它必含 $\mathfrak{n}^{+}$-不变向量，这与 7.15.5(i) 矛盾）。故由 $\nu$ 定义的从 $\Gamma\delta$ 到余自由 $N^{+}$-模的态射是单射。再由维数理由，它是同构。}。于是 $H^{\cdot}(\mathfrak{n}^{+}_x,\Gamma\delta)=H^{0}(\mathfrak{n}^{+}_x,\Gamma\delta)^{\chi}=\mathbb{C}\cdot vac$。又 $H^{\cdot}i_x^{!}\delta=H^{0}i_x^{!}\delta=\mathbb{C}\cdot vac$，故得所要的同构。
